% Chapter 4 of The Replica Rosetta Stone (The Missing Textbooks, volume 7):
% Bayes-optimal inference and the Nishimori line.
\rstchapter{Bayes-Optimal Inference}{Why replica symmetry holds}{the Nishimori line\dict overlap concentration}{ch:nishimori}
\stonefoot{proved}
\begin{refsection}[chapters/04-nishimori.bib]

\begin{chaptersummary}
When the Gibbs measure is the true Bayesian posterior, the hidden signal behaves like one more sample from it, and the replica-symmetric answer of physics is exact. We derive the mutual information and the minimum mean-square error of $\Z_2$ synchronization by replicas in which the signal is replica number zero, and prove the formula completely: Guerra's interpolation bounds the mutual information at every size, a small side channel makes the overlap concentrate, and an adaptive interpolation turns the concentration into the matching bound. Exact enumeration at $N\le12$ and message passing at $N=2000$ confirm the curve. Replica symmetry on the Nishimori line beyond dense models with independent priors remains open.
\end{chaptersummary}

\section{The question: signs from noisy pairwise products}\label{sec:nishimori:question}

Each of $N$ objects carries a hidden sign $x^*_i\in\{-1,+1\}$, and for every pair we observe the product $x^*_ix^*_j$ through Gaussian noise. How much can be learned about the signs, up to the global flip that no pairwise measurement can see? This is \term{$\Z_2$ synchronization}\index{Z2 synchronization@$\Z_2$ synchronization},\sidenote{In random-matrix language it is the spiked Wigner model with the Rademacher prior; \cref{ch:spiked} studies its spectrum and \cref{ch:amp} its algorithms.} and in the book's normalization the data are
\begin{equation}\label{eq:nishimori:model}
A=\frac\lambda N\,x^*x^{*\top}+W,
\end{equation}
with $x^*$ uniform on $\Sigma_N=\{-1,+1\}^N$ and, independent of it, $W\sim\mathrm{GOE}(N)$: symmetric, with independent centered Gaussian entries of variance $1/N$ off the diagonal and $2/N$ on it, so that its density is proportional to $e^{-N\tr(W^2)/4}$. This is the \term{spiked Wigner model}\index{spiked Wigner model} with the Rademacher prior; its top eigenvalue leaves the bulk when $\lambda>1$ (\cref{thm:spiked:bbp}). We ask what the best estimator can do.

\paragraph{The posterior is a spin glass.} The likelihood of $A$ given $x^*=x$ is proportional to $\exp(-\frac N4\norm{A-\frac\lambda Nxx^\top}_F^2)$, and $\norm x^4=N^2$ on $\Sigma_N$, so the posterior is the Gibbs measure
\begin{equation}\label{eq:nishimori:posterior}
\Prob(x^*=x\mid A)=\frac{e^{H_N(x)}}{\sum_{x'\in\Sigma_N}e^{H_N(x')}},\qquad H_N(x)=\frac\lambda2\,x^\top Ax.
\end{equation}
We write $\gibbs f$ for the posterior mean of a function of independent posterior samples $x^1,x^2,\dots$, the \term{replicas}\index{replica}, $\E$ for the expectation over $(x^*,W)$, and
\[
R_{12}=\frac1N\ip{x^1}{x^2},\qquad R_{1*}=\frac1N\ip{x^1}{x^*}
\]
for the overlap of two replicas and of a replica with the truth, which physicists call the magnetization $m$. In the variables $\sigma_i=x^*_ix_i$, and up to a constant, $H_N$ is the sum of $\frac\lambda{\sqrt N}\sum_{i<j}J_{ij}\sigma_i\sigma_j$ and $\frac{\lambda^2}N\sum_{i<j}\sigma_i\sigma_j$, where the $J_{ij}=x^*_ix^*_j\sqrt NW_{ij}$ are again independent standard Gaussians: the Sherrington--Kirkpatrick model at inverse temperature $\beta=\lambda$ (\cref{ch:sk}) with couplings of mean $J_0/N$ and variance $J^2/N$, $\beta J=\lambda$, $\beta J_0=\lambda^2$. The relation $\beta=J_0/J^2$ defines the \term{Nishimori line}\index{Nishimori line}, on which \textcite{Nishimori1981} found exact results by a gauge transformation; \textcite{Iba1999} pointed out that on it the Gibbs measure is a Bayesian posterior.\sidenote{Nishimori worked with the $\pm J$ model, in which each coupling is $+J$ with probability $p$ and $-J$ otherwise; there the line is $\beta J=\frac12\log\frac p{1-p}$.}

\paragraph{What we want to know.} The \term{evidence}\index{evidence} is the likelihood ratio of the spiked model against pure noise,
\[
\cZ_N=\frac{p_\lambda(A)}{p_0(A)}=2^{-N}\sum_{x\in\Sigma_N}e^{H_N(x)-\lambda^2N/4},
\]
and $\phi_N=\frac1N\E\log\cZ_N$ is the \term{free entropy}\index{free entropy}. Since $\log\frac{p_\lambda(A\mid x^*)}{p_0(A)}=H_N(x^*)-\frac{\lambda^2N}4$ has mean $\frac{\lambda^2N}4$, the \term{mutual information}\index{mutual information} is
\begin{equation}\label{eq:nishimori:mi}
\frac1NI(x^*;A)=\frac{\lambda^2}4-\phi_N.
\end{equation}
The sign of $x^*$ is not identifiable, so we estimate $x^*x^{*\top}$. The posterior mean is the best estimator in mean square, and its error is the \term{minimum mean-square error}\index{minimum mean-square error}
\begin{equation}\label{eq:nishimori:mmse}
\mathrm{MMSE}_N=\frac1{N^2}\E\bigl\|x^*x^{*\top}-\gibbs{xx^\top}\bigr\|_F^2=1-\E\gibbs{R_{1*}^2},
\end{equation}
by expanding the square and the Nishimori identity $\E\gibbs{R_{12}^2}=\E\gibbs{R_{1*}^2}$ of \cref{lem:nishimori:nishimori}; the estimator $0$ has error $1$. The prediction involves the \term{scalar channel}\index{scalar channel} $Y=\sqrt\gamma X+Z$, with $X$ uniform on $\{\pm1\}$ and $Z\sim\Normal(0,1)$ independent, whose posterior mean is $\tanh(\sqrt\gamma Y)$. Its overlap $\E[X\tanh(\sqrt\gamma Y)]$ and its free entropy (\cref{ch:freeentropy}) are
\begin{equation}\label{eq:nishimori:scalar}
G(\gamma)=\E\tanh(\gamma+\sqrt\gamma Z)=1-\mathrm{mmse}(\gamma),\qquad\psi(\gamma)=\E\log\cosh(\gamma+\sqrt\gamma Z)-\frac\gamma2,
\end{equation}
related by the I-MMSE identity $\psi'=\frac12G$ (\cref{thm:freeentropy:immse}).

\begin{physicsbox}[Lesieur, Krzakala and Zdeborová 2017]
For $\Z_2$ synchronization, as $N\to\infty$,
\begin{equation}\label{eq:nishimori:prediction}
\frac1NI(x^*;A)\to\frac{\lambda^2}4-\max_{0\le q\le1}\phi_{\mathrm{RS}}(\lambda,q),\qquad\phi_{\mathrm{RS}}(\lambda,q)=\psi(\lambda^2q)-\frac{\lambda^2q^2}4,
\end{equation}
and $\mathrm{MMSE}_N\to1-q_*^2$, where the maximizer $q_*=q_*(\lambda)$ solves
\begin{equation}\label{eq:nishimori:fixedpoint}
q=\E\tanh\bigl(\lambda^2q+\lambda\sqrt q\,Z\bigr).
\end{equation}
For $\lambda\le1$ the only solution is $q_*=0$ and nothing can be learned; for $\lambda>1$ there is a positive one, with $q_*\approx2(\lambda-1)$ near the threshold. The replica-symmetric ansatz behind the formula is exact, because the posterior sits on the Nishimori line.
\end{physicsbox}

How do physicists derive this, and why do they trust replica symmetry here when it fails for the Sherrington--Kirkpatrick model at low temperature? And what is proved? Everything in the box (\cref{thm:nishimori:z2}), for the reason the physicists gave.

\section{The physics: the truth is a replica}\label{sec:nishimori:physics}

The replica method computes $\E\log\cZ_N$ from the moments $\E\cZ_N^n$ (\cref{ch:freeentropy}). On the Nishimori line the signal enters them as one more replica.

\begin{computation}{The replicated evidence}\label{comp:nishimori:replicated}
\step \emph{Replicate and average the noise.} $\cZ_N^n$ is the mean, over independent uniform $x^1,\dots,x^n\in\Sigma_N$, of $\exp\sum_a(H_N(x^a)-\frac{\lambda^2N}4)$, and $H_N(x^a)=\frac{\lambda^2N}2R_{0a}^2+\frac\lambda2x^{a\top}Wx^a$, with $x^0=x^*$ and $R_{ab}=\frac1N\ip{x^a}{x^b}$. Given the replicas, $\frac\lambda2\sum_ax^{a\top}Wx^a$ is a centered Gaussian, and the covariance $\E W_{ij}W_{kl}=\frac1N(\delta_{ik}\delta_{jl}+\delta_{il}\delta_{jk})$ of the GOE makes its variance $\frac{\lambda^2}{2N}\sum_{a,b}\ip{x^a}{x^b}^2=\frac{\lambda^2Nn}2+\lambda^2N\sum_{1\le a<b}R_{ab}^2$.
\step \emph{Collect.} Half the variance, $\frac{\lambda^2Nn}4+\frac{\lambda^2N}2\sum_{1\le a<b}R_{ab}^2$, appears in the exponent, its first term cancels the normalization, and with $x^0,\dots,x^n$ independent and uniform on $\Sigma_N$, for every $N$ and every integer $n\ge1$,
\begin{equation}\label{eq:nishimori:replicated}
\E\cZ_N^n=\E\exp\Bigl(\frac{\lambda^2N}2\sum_{0\le a<b\le n}R_{ab}^2\Bigr).
\end{equation}
\step \emph{Read.} The weight is symmetric under all permutations of $x^0,\dots,x^n$: the truth enters as replica number zero. At $n=1$, $\E\cZ_N$ is the second moment of the likelihood ratio under pure noise (\cref{prop:nishimori:below}).
\end{computation}

\begin{heuristic}[Physical derivation: replicas on the Nishimori line]
Decouple each pair by a Gaussian integral,
\[
e^{\lambda^2NR^2/2}=\sqrt{\lambda^2N/2\pi}\int_\R e^{-\lambda^2N(Q^2/2-QR)}\dd Q.
\]
The sum over the replicas then factorizes over sites. With $0\le a<b\le n$, $\sigma$ uniform on $\{\pm1\}^{n+1}$ and $C_{N,n}$ subexponential in $N$, this gives $\E\cZ_N^n=C_{N,n}\int e^{NS_n(Q)}\prod_{a<b}\dd Q_{ab}$ with
\[
S_n(Q)=-\frac{\lambda^2}2\sum_{a<b}Q_{ab}^2+\log\E_\sigma e^{\lambda^2\sum_{a<b}Q_{ab}\sigma_a\sigma_b}.
\]
At a stationary point $Q_{ab}$ is the overlap $\langle\sigma_a\sigma_b\rangle$ of replicas $a$ and $b$ in the single-site replicated measure. The replica-symmetric ansatz $Q_{ab}=q$ for all pairs treats the truth as one more replica, as \eqref{eq:nishimori:replicated} invites. Since $\sum_{a<b}\sigma_a\sigma_b=\frac12(\sum_a\sigma_a)^2-\frac{n+1}2$, a second Gaussian integral gives
\[
S_n(q)=-\binom{n+1}2\frac{\lambda^2q^2}2-\frac{(n+1)\lambda^2q}2+\log\E\cosh^{n+1}(\lambda\sqrt q\,Z).
\]
Now continue in $n$. The function $S_0$ vanishes, as it must since $\E\cZ_N^0=1$, and exchanging $n\to0$ with $N\to\infty$ in $\frac1N\E\log\cZ_N=\lim_{n\to0}\frac1{nN}\log\E\cZ_N^n$ leaves $\partial_nS_n(q)$ at $n=0$. With $\gamma=\lambda^2q$ and the shift $e^{-\gamma/2}\E[e^{\sqrt\gamma Z}f(\sqrt\gamma Z)]=\E f(\gamma+\sqrt\gamma Z)$ for even $f$,
\[
\partial_nS_n(q)\big|_{n=0}=-\frac{\lambda^2q^2}4-\frac\gamma2+\frac{\E[\cosh(\sqrt\gamma Z)\log\cosh(\sqrt\gamma Z)]}{\E\cosh(\sqrt\gamma Z)}=\psi(\lambda^2q)-\frac{\lambda^2q^2}4,
\]
which is $\phi_{\mathrm{RS}}(\lambda,q)$; stationarity in $q$ is \eqref{eq:nishimori:fixedpoint}. Which stationary point? For integer $n$ Laplace's method takes the maximum of $S_n$. Near $n=0$, $S_n\approx n\,\phi_{\mathrm{RS}}$ with $n>0$, and the number $\binom{n+1}2$ of order parameters stays positive, so the maximum survives: $\lim\phi_N=\max_q\phi_{\mathrm{RS}}(\lambda,q)$, which is \eqref{eq:nishimori:prediction} by \eqref{eq:nishimori:mi}. (In the Sherrington--Kirkpatrick model the replicas are $x^1,\dots,x^n$ alone, $\binom n2<0$, and the rule becomes a minimum; \cref{ch:sk}.)
\end{heuristic}

\paragraph{Why replica symmetry.} The guess is wrong below the de Almeida--Thouless line of the Sherrington--Kirkpatrick model; on the Nishimori line physicists trust it for three reasons. The weight \eqref{eq:nishimori:replicated} is symmetric in all $n+1$ replicas, so a symmetric saddle point has $Q_{0a}=Q_{ab}$. At finite $N$ the overlap of two posterior samples has the law of the overlap of one sample with the truth, $P(q)=P(m)$, which \textcite{NishimoriSherrington2001} proved for the $\pm J$ model in every dimension. And the magnetization $m=R_{1*}$ self-averages, being the derivative of a self-averaging free energy in its conjugate field; so $P(m)$, and with it $P(q)$, has no spread beyond the atoms $\pm m$ forced by the global flip. \Textcite[\S II.G]{ZdeborovaKrzakala2016} note that this last step is open in general.

\paragraph{What the derivation assumed.} Five steps are taken on trust: the continuation from integer $n$ to $n\to0$, which the integer moments do not determine in general (\cref{ch:freeentropy}); the exchange of $n\to0$ with $N\to\infty$; Laplace's method over a number $\binom{n+1}2$ of variables that is not an integer; the replica-symmetric ansatz; and the choice of the global maximum.

\paragraph{Continuous and discontinuous transitions.} The fixed-point equation is $q=G(\lambda^2q)$, and $G(\gamma)=\gamma-\gamma^2+\frac53\gamma^3-\cdots$ (\cref{ex:nishimori:series}). At $\lambda=1$ a positive solution bifurcates from $0$, with $q_*=2\epsilon-\frac13\epsilon^2+O(\epsilon^3)$ at $\lambda=1+\epsilon$: the transition is continuous and sits at the spectral threshold (\cref{fig:nishimori:potential}(a)). For a prior $P_0$ with $\E X^2=1$ the derivation is the same, with $\psi$ replaced by $\psi_{P_0}(\gamma)=\E\log\int e^{\sqrt\gamma Zx+\gamma xX-\gamma x^2/2}P_0(\dd x)$, $X\sim P_0$. For the sparse Rademacher prior, with mass $1-\rho$ at $0$ and $\frac\rho2$ at $\pm\rho^{-1/2}$, a second local maximum of $\phi_{\mathrm{RS}}$ appears while $q=0$ is still a local maximum, and overtakes it at $\lambda_{\mathrm{IT}}$, which is $0.8645$ for $\rho=0.05$ (\cref{fig:nishimori:potential}(b)): the optimal overlap jumps, a first-order transition, and up to a value slightly above $1$ an algorithm that climbs the potential from $q=0$ stays near it, the hard phase (\cref{ch:spiked}).

\begin{figure}[t]
\begin{fullwidth}
\centering
\begin{tikzpicture}
\begin{groupplot}[group style={group size=2 by 1,horizontal sep=17mm},rplot,width=80mm,height=45mm]
\nextgroupplot[xlabel={overlap $q$},ylabel={$\phi_{\mathrm{RS}}(\lambda,q)$},xmin=0,xmax=1,ymin=-0.095,ymax=0.08,
  title={(a) Rademacher prior},legend pos=south west,legend style={font=\scriptsize}]
\addplot[physics!30,line width=0.9pt,no marks] table[x=q,y=r08] {figures/data/04-nishimori-potential.dat};
\addplot[physics!55,line width=0.9pt,no marks] table[x=q,y=r10] {figures/data/04-nishimori-potential.dat};
\addplot[physics!78,line width=0.9pt,no marks] table[x=q,y=r12] {figures/data/04-nishimori-potential.dat};
\addplot[physics,line width=1.0pt,no marks] table[x=q,y=r15] {figures/data/04-nishimori-potential.dat};
\addplot[guide,forget plot] coordinates {(0,0) (1,0)};
\legend{$\lambda=0.8$,$\lambda=1$,$\lambda=1.2$,$\lambda=1.5$}
\nextgroupplot[xlabel={overlap $q$},xmin=0,xmax=1,ymin=-0.0042,ymax=0.0042,
  title={(b) sparse Rademacher prior, $\rho=0.05$},legend pos=south west,legend style={font=\scriptsize},
  yticklabel style={/pgf/number format/fixed,/pgf/number format/precision=3}]
\addplot[physics!30,line width=0.9pt,no marks] table[x=q,y=s084] {figures/data/04-nishimori-potential.dat};
\addplot[physics!65,line width=0.9pt,no marks] table[x=q,y=sIT] {figures/data/04-nishimori-potential.dat};
\addplot[physics,line width=1.0pt,no marks] table[x=q,y=s088] {figures/data/04-nishimori-potential.dat};
\addplot[guide,forget plot] coordinates {(0,0) (1,0)};
\legend{$\lambda=0.84$,$\lambda=\lambda_{\mathrm{IT}}=0.8645$,$\lambda=0.88$}
\end{groupplot}
\end{tikzpicture}
\end{fullwidth}
\caption{The replica-symmetric potential $\phi_{\mathrm{RS}}(\lambda,q)$. (a) Rademacher prior, $\lambda=0.8$, $1$, $1.2$, $1.5$ (light to dark): the maximizer leaves $q=0$ continuously at $\lambda=1$ ($q_*=0.358$ at $\lambda=1.2$, $0.692$ at $1.5$). (b) Sparse Rademacher prior, $\rho=0.05$, $\lambda=0.84$, $\lambda_{\mathrm{IT}}=0.8645$, $0.88$: a second maximum appears, ties with $q=0$ at $\lambda_{\mathrm{IT}}$ (at $q=0.585$), and then dominates. Gauss--Hermite quadrature, $201$ nodes; script \texttt{04-nishimori-rs.py}.}\label{fig:nishimori:potential}
\end{figure}

\section{The stone}\label{sec:nishimori:stone}

The proof in \cref{sec:nishimori:proofguerra,sec:nishimori:proofmatching} follows the derivation step by step, except at the continuation in $n$, which it replaces by two inequalities.

\begin{stone}{replicas on the Nishimori line}{Guerra's interpolation with a side channel}
\stonestep{Set up}{Average $n$ copies of the evidence; the signal enters as replica zero (\cref{comp:nishimori:replicated}).}{Interpolate between the matrix channel and $N$ scalar channels $y=\sqrt r\,x^*+z$: the free entropy $\phi_N(t,r)$ (\cref{sec:nishimori:interp}).}
\stonestep{The Nishimori symmetry}{The replicated weight is symmetric in $x^0,\dots,x^n$: the truth is statistically a replica, and $P(q)=P(m)$.}{$\E\gibbs{f(x^1,\dots,x^k,x^*)}=\E\gibbs{f(x^1,\dots,x^{k+1})}$ at every $N$, by Bayes' rule (\cref{lem:nishimori:nishimori}); for signs, also the gauge (\cref{lem:nishimori:gauge}).}
\stonestep{Order parameters}{Gaussian integrals decouple the sites; overlaps $Q_{ab}$ and conjugate fields $\lambda^2Q_{ab}$ meet at a saddle point.}{Integration by parts and the Nishimori identity: $\partial_t\phi_N=\frac{\lambda^2}4\E\gibbs{R_{1*}^2}$, $\partial_r\phi_N=\frac12\E\gibbs{R_{1*}}$; $r$ is the conjugate field (\cref{lem:nishimori:derivatives}).}
\stonestep{Replica symmetry, and a maximum}{The ansatz $Q_{ab}=q$ for all pairs, the truth included; $\binom{n+1}2>0$ keeps the saddle point a maximum ($\binom n2<0$ makes it a minimum in the Sherrington--Kirkpatrick model).}{No ansatz: the remainder $\frac{\lambda^2}4\int_0^1\E\gibbs{(R_{1*}-q)^2}\dd t$ is a variance, so $\phi_N\ge\phi_{\mathrm{RS}}(\lambda,q)$ for every $q$ and $N$; for Sherrington--Kirkpatrick it has the other sign (\cref{sec:nishimori:guerra}).}
\stonestep{Self-averaging of $m$}{$P(m)$ is trivial because $m$ is the derivative of a self-averaging free energy in its conjugate field.}{Convexity in $r$ and Gaussian concentration make $\E\gibbs{(R_{1*}-\E\gibbs{R_{1*}})^2}$ small for most $r$ (\cref{prop:nishimori:overlap}).}
\stonestep{The saddle point is the answer}{Continue to $n\to0$, exchange the limits, keep the global maximum.}{No counterpart for the continuation. The adaptive path $q(t)=\E\gibbs{R_{1*}}_{t,r(t)}$ makes the remainder the overlap variance, small on average over $r$ (\cref{sec:nishimori:matching}).}
\stonestep{The fixed point}{$q=\E\tanh(\lambda^2q+\lambda\sqrt qZ)$, continuous at $\lambda=1$ for signs; first order for sparse priors.}{The binary mmse is completely monotone, so the maximizer is unique (\cref{lem:nishimori:cm}); for other priors the formula is a theorem (\cref{thm:nishimori:rsformula}).}
\end{stone}

The Nishimori identity does in the proof what the permutation symmetry of the $n+1$ replicas does in the derivation, and the physicists' reason for replica symmetry, a magnetization conjugate to a field, becomes a theorem once the field is made real, as a side channel, and averaged over. What the proof cannot keep is the single saddle point. In its place stand two inequalities of different natures: Guerra's bound holds at every $N$ and needs no concentration, while the matching bound holds only in the limit and needs both the concentration and the side channel.

\section{The proof: an interpolation with a sign}\label{sec:nishimori:proofguerra}

The general theorem is due to Lelarge and Miolane; we state it in the book's normalization ($\mathrm{snr}=\lambda^2$) and then prove the case of signs completely, in this section and the next.

\begin{theorem}[\textcite{LelargeMiolane2019}, Theorem~1]\label{thm:nishimori:rsformula}
Let $P_0$ be a probability measure on $\R$ with $v=\int x^2P_0(\dd x)<\infty$, let $x^*_1,\dots,x^*_N$ be independent with law $P_0$, and observe $Y_{ij}=\frac\lambda{\sqrt N}x^*_ix^*_j+Z_{ij}$ for $1\le i<j\le N$, with independent standard Gaussians $Z_{ij}$, independent of $x^*$. With $\phi_{\mathrm{RS}}(\lambda,q)=\psi_{P_0}(\lambda^2q)-\frac{\lambda^2q^2}4$, for every $\lambda>0$,
\[
\lim_{N\to\infty}\frac1NI(x^*;Y)=\frac{\lambda^2v^2}4-\sup_{q\ge0}\phi_{\mathrm{RS}}(\lambda,q).
\]
The limit is concave in $\lambda^2$. At each point where it is differentiable, which is every $\lambda>0$ outside a countable set, the maximizer $q_*(\lambda)$ is unique and
\[
\lim_{N\to\infty}\frac2{N(N-1)}\sum_{i<j}\E\bigl(x^*_ix^*_j-\E[x^*_ix^*_j\mid Y]\bigr)^2=v^2-q_*(\lambda)^2.
\]
\end{theorem}

The $Y_{ij}$ are the entries of $\sqrt NA$ above the diagonal; Lelarge and Miolane remark that the case in which all entries are observed can be deduced from theirs. For signs the diagonal $A_{ii}=\frac\lambda N+W_{ii}$ is independent of everything else, so $I(x^*;A)=I(x^*;Y)$, and $\mathrm{MMSE}_N$ is $\frac{N-1}N$ times their error.

\begin{theorem}[The replica-symmetric formula for $\Z_2$ synchronization]\label{thm:nishimori:z2}
Let $x^*$ be uniform on $\Sigma_N$ in \eqref{eq:nishimori:model}, and let $\lambda\ge0$.
\begin{enumerate}
\item For every $N\ge1$ and every $q\ge0$, $\phi_N\ge\phi_{\mathrm{RS}}(\lambda,q)$; equivalently $\frac1NI(x^*;A)\le\frac{\lambda^2}4-\phi_{\mathrm{RS}}(\lambda,q)$.
\item There is $C(\lambda)<\infty$ with $\phi_N\le\max_{0\le q\le1}\phi_{\mathrm{RS}}(\lambda,q)+C(\lambda)N^{-1/20}$ for $N\ge11$; so $\phi_N\to\max_q\phi_{\mathrm{RS}}(\lambda,q)$.
\item The maximizer $q_*(\lambda)$ is unique: $q_*=0$ for $\lambda\le1$, and for $\lambda>1$ it is the unique solution of \eqref{eq:nishimori:fixedpoint} in $(0,1]$.
\item $\mathrm{MMSE}_N\to1-q_*(\lambda)^2$ for every $\lambda\ge0$.
\end{enumerate}
\end{theorem}

Parts 1 and 2 give \cref{thm:nishimori:rsformula} for this prior, which \textcite{DeshpandeAbbeMontanari2016} had proved by message passing and interpolation; part 3 removes the exceptional set from part 4. We prove part 1 here, part 2 in \cref{sec:nishimori:matching} and parts 3 and 4 in \cref{sec:nishimori:mmse}.

\subsection{The interpolating model and the Nishimori identity}\label{sec:nishimori:interp}

For $t\in[0,1]$ and $r\ge0$ observe
\begin{equation}\label{eq:nishimori:interp}
A_t=\frac{\lambda\sqrt t}N\,x^*x^{*\top}+W,\qquad y_r=\sqrt r\,x^*+z,
\end{equation}
with $z\sim\Normal(0,I_N)$ independent of $(x^*,W)$: the matrix channel at signal-to-noise ratio $\lambda\sqrt t$ and $N$ scalar channels at ratio $r$. On $\Sigma_N$ the terms $\frac{\lambda^2t}{4N}\norm x^4$ and $\frac r2\norm x^2$ of the log-likelihood are constant, so the posterior is the Gibbs measure $\gibbs\cdot_{t,r}$ with Hamiltonian $H_{t,r}(x)=\frac{\lambda\sqrt t}2x^\top A_tx+\sqrt r\ip{y_r}x$, the evidence is $\cZ_{t,r}=2^{-N}\sum_x\exp(H_{t,r}(x)-\frac{\lambda^2tN}4-\frac{rN}2)$, and we write $\phi_N(t,r)=\frac1N\E\log\cZ_{t,r}$. At $(t,r)=(1,0)$ this is the model. At $t=0$ the Gibbs measure is a product, $\cZ_{0,r}=\prod_ie^{-r/2}\cosh(\sqrt r(y_r)_i)$, and $\phi_N(0,r)=\psi(r)$.

\begin{lemma}[Nishimori identity]\label{lem:nishimori:nishimori}
For every $t\in[0,1]$, $r\ge0$, $k\ge1$ and $f\colon\Sigma_N^{k+1}\to\R$,
$\E\gibbs{f(x^1,\dots,x^k,x^*)}_{t,r}=\E\gibbs{f(x^1,\dots,x^k,x^{k+1})}_{t,r}$.
\end{lemma}

\begin{proof}
Write $D=(A_t,y_r)$ and $\pi(\cdot\mid D)$ for the posterior. By Bayes' rule $x^*$ has the conditional law $\pi(\cdot\mid D)$ given $D$, so, writing $x^{k+1}$ for $x^*$,
\[
\E\gibbs{f(x^1,\dots,x^k,x^*)}=\E_D\sum_{x^1,\dots,x^{k+1}}f(x^1,\dots,x^{k+1})\prod_{a=1}^{k+1}\pi(x^a\mid D),
\]
which is $\E\gibbs{f(x^1,\dots,x^{k+1})}$.\qedhere
\end{proof}

This is the Bayesian form of Nishimori's identity (\cref{lem:freeentropy:nishimori}; \textcite[Proposition~16]{LelargeMiolane2019}).\sidenote{\Textcite{BarbierMacris2019} remark that the name is used loosely for this consequence of Bayes' rule, and that Nishimori's identity proper uses a gauge symmetry to remove the signal. For signs we have both.} We use it as $\E\gibbs{R_{1*}}=\E\gibbs{R_{12}}=\frac1N\E\norm{\gibbs x}^2\in[0,1]$, $\E\gibbs{R_{1*}^2}=\E\gibbs{R_{12}^2}$ and $\E\gibbs{R_{1*}R_{12}}=\E\gibbs{R_{13}R_{12}}$.

\begin{lemma}[Gauge symmetry]\label{lem:nishimori:gauge}
Given $x^*=s\in\Sigma_N$, the law of $\cZ_{t,r}$ does not depend on $s$; so $\log\cZ_{t,r}$ is independent of $x^*$.
\end{lemma}

\begin{proof}
Let $D_s=\diag(s)$. Given $x^*=s$, the entries $s_is_jW_{ij}$ of $D_sWD_s$ are independent with the variances of those of $W$, so $(D_sWD_s,D_sz)$ has the law of $(W,z)$. The substitution $x=D_s\sigma$ is a bijection of $\Sigma_N$, and since $D_sx^*=\mathbf1$ it turns $H_{t,r}$ into the Hamiltonian of the model with truth $\mathbf1$ and noise $(D_sWD_s,D_sz)$.
\end{proof}

\subsection{Two derivatives and a sum rule}

\begin{lemma}[Derivatives of the free entropy]\label{lem:nishimori:derivatives}
The function $\phi_N$ is continuously differentiable on $[0,1]\times[0,\infty)$, with
\[
\partial_t\phi_N(t,r)=\frac{\lambda^2}4\E\gibbs{R_{1*}^2}_{t,r},\qquad\partial_r\phi_N(t,r)=\frac12\E\gibbs{R_{1*}}_{t,r}.
\]
Moreover $\E\gibbs{R_{1*}}_{t,r}$ is nondecreasing in $r$, and $\E\gibbs{R_{1*}^2}_{t,r}$ is nondecreasing in $t$.
\end{lemma}

\begin{proof}
The plan is to differentiate under the expectation, to remove the noise by Gaussian integration by parts (\cref{thm:toolkit:gaussianibp}), and to remove the replicas by the Nishimori identity. For fixed $N$ all quantities are smooth for $t,r>0$ and dominated by polynomials in $(W,z)$; the formulas extend to $t=0$ and $r=0$ by continuity of their right sides.

\emph{In $r$.} Since $\sqrt r\ip{y_r}x=r\ip{x^*}x+\sqrt r\ip zx$, $\partial_r\log\cZ_{t,r}=\gibbs{\ip{x^*}x}+\frac1{2\sqrt r}\gibbs{\ip zx}-\frac N2$. Integration by parts in $z_i$, with $\partial_{z_i}\gibbs{x_i}=\sqrt r(1-\gibbs{x_i}^2)$, gives $\E\gibbs{\ip zx}=\sqrt r\,\E(N-\norm{\gibbs x}^2)$, so $\partial_r\E\log\cZ_{t,r}=N\E\gibbs{R_{1*}}-\frac N2\E\gibbs{R_{12}}=\frac N2\E\gibbs{R_{1*}}$.

\emph{In $t$.} Here $\frac{\lambda\sqrt t}2x^\top A_tx=\frac{\lambda^2t}{2N}\ip{x^*}x^2+\frac{\lambda\sqrt t}2x^\top Wx$, and $H_{t,r}$ depends on $W_{ij}$, $i<j$, through $\lambda\sqrt tW_{ij}x_ix_j$ and on $W_{ii}$ only through a constant. Integration by parts with $\E W_{ij}^2=\frac1N$ gives $\E\gibbs{x^\top Wx}=\frac{2\lambda\sqrt t}N\sum_{i<j}\E(1-\gibbs{x_ix_j}^2)=\lambda\sqrt tN(1-\E\gibbs{R_{12}^2})$, so that
\[
\partial_t\E\log\cZ_{t,r}=\frac{\lambda^2N}2\E\gibbs{R_{1*}^2}+\frac{\lambda^2N}4\bigl(1-\E\gibbs{R_{12}^2}\bigr)-\frac{\lambda^2N}4=\frac{\lambda^2N}4\E\gibbs{R_{1*}^2}.
\]
Both derivatives are continuous in $(t,r)$ by dominated convergence.

\emph{Monotonicity.} By the tower property $\E\gibbs{R_{1*}}_{t,r}=1-\frac1N\E\norm{x^*-\E[x^*\mid A_t,y_r]}^2$. For $r'<r$ and an independent $\xi\sim\Normal(0,I_N)$, the vector $\sqrt{r'/r}\,y_r+\sqrt{1-r'/r}\,\xi$ has, jointly with $(x^*,A_t)$, the law of $y_{r'}$, and conditioning on more information lowers the mean-square error. The same argument with $\sqrt{t'/t}\,A_t+\sqrt{1-t'/t}\,W'$, for an independent GOE matrix $W'$, and \eqref{eq:nishimori:mmse} for the interpolating model gives the monotonicity in $t$.
\end{proof}

\begin{proposition}[The sum rule]\label{prop:nishimori:interpolation}
Let $q\colon[0,1]\to[0,\infty)$ be continuous, $\epsilon\ge0$, and $r(t)=\epsilon+\lambda^2\int_t^1q(u)\dd u$. Then
\begin{multline*}
\phi_N(1,\epsilon)=\psi\Bigl(\epsilon+\lambda^2\int_0^1q(t)\dd t\Bigr)-\frac{\lambda^2}4\int_0^1q(t)^2\dd t\\
+\frac{\lambda^2}4\int_0^1\E\gibbs{(R_{1*}-q(t))^2}_{t,r(t)}\dd t.
\end{multline*}
\end{proposition}

\begin{proof}
Let $\varphi(t)=\phi_N(t,r(t))$, so that $\varphi(1)=\phi_N(1,\epsilon)$ and $\varphi(0)=\psi(r(0))$. Since $r'=-\lambda^2q$, the chain rule and \cref{lem:nishimori:derivatives} give, with the brackets at $(t,r(t))$,
\[
\varphi'(t)=\frac{\lambda^2}4\E\gibbs{R_{1*}^2}-\frac{\lambda^2q(t)}2\E\gibbs{R_{1*}}=\frac{\lambda^2}4\E\gibbs{(R_{1*}-q(t))^2}-\frac{\lambda^2q(t)^2}4,
\]
and integrating over $[0,1]$ gives the claim.
\end{proof}

\subsection{Guerra's bound}\label{sec:nishimori:guerra}

\begin{proof}[Proof of part 1 of \cref{thm:nishimori:z2}]
Take $\epsilon=0$ and a constant $q\ge0$ in \cref{prop:nishimori:interpolation}. The first two terms are $\phi_{\mathrm{RS}}(\lambda,q)$, and the last is nonnegative:
\[
\phi_N-\phi_{\mathrm{RS}}(\lambda,q)=\frac{\lambda^2}4\int_0^1\E\gibbs{(R_{1*}-q)^2}_{t,\lambda^2q(1-t)}\dd t\ge0.
\]
The second form follows from \eqref{eq:nishimori:mi}. Since $\partial_q\phi_{\mathrm{RS}}=\frac{\lambda^2}2(G(\lambda^2q)-q)<0$ for $q\ge1$, the supremum over $q\ge0$ is a maximum over $[0,1]$.\qedhere
\end{proof}

The identity holds at every $N$ with no error term, because the variance $2/N$ of the diagonal of the GOE makes the terms in $\norm x^4$ cancel; \textcite{KrzakalaXuZdeborova2016} proved the bound at every $N$ for priors with finite support. It needs no concentration and no perturbation: it is a sign. For the Sherrington--Kirkpatrick model the same interpolation has derivative $\frac{\beta^2}4((1-q)^2-\E\gibbs{(R_{12}-q)^2})$, and Guerra's bound is an upper bound (\cref{thm:sk:rsbound}; \cref{ex:nishimori:sk}). Here the planted term contributes $\frac{\lambda^2}2\E\gibbs{R_{1*}^2}$, twice what the noise removes, and the Nishimori identity turns the balance into $+\frac{\lambda^2}4\E\gibbs{(R_{1*}-q)^2}$. The physicist's rule, a maximum here and a minimum there, is this sign. Below the threshold the bound has a partner, from the second moment of \cref{comp:nishimori:replicated}.

\begin{proposition}[Below the threshold]\label{prop:nishimori:below}
For $0\le\lambda<1$ and every $N\ge1$, $\displaystyle0\le\phi_N\le\frac{-\log(1-\lambda^2)}{2N}$.
\end{proposition}

\begin{proof}
The lower bound is part 1 at $q=0$. For the upper bound, Jensen's inequality and \eqref{eq:nishimori:replicated} at $n=1$ give $N\phi_N\le\log\E\cZ_N=\log\E e^{\lambda^2NR_{01}^2/2}$. Here $\sqrt NR_{01}=N^{-1/2}\sum_i\sigma_i$ with independent uniform signs $\sigma_i=x^0_ix^1_i$, so $\E e^{\theta\sqrt NR_{01}}=\cosh(\theta/\sqrt N)^N\le e^{\theta^2/2}$, and with an independent $\xi\sim\Normal(0,1)$,
\[
\E e^{\lambda^2NR_{01}^2/2}=\E e^{\lambda\xi\sqrt NR_{01}}\le\E e^{\lambda^2\xi^2/2}=(1-\lambda^2)^{-1/2}.\qedhere
\]
\end{proof}

\section{The proof: overlap concentration and the matching bound}\label{sec:nishimori:proofmatching}

Guerra's remainder would vanish if the overlap concentrated along the path, which is what replica symmetry claims. The field conjugate to the overlap is the ratio $r$ of the side channel, $\partial_r\phi_N=\frac12\E\gibbs{R_{1*}}$, and a free entropy that concentrates and is convex in $r$ has a derivative that concentrates for most $r$: the physicists' argument, whose deterministic core is the next lemma.\index{overlap!concentration}

\begin{lemma}[Derivatives of close convex functions]\label{lem:nishimori:convex}
Let $f$ and $g$ be convex and differentiable on an open interval that contains $[r-\delta,r+\delta]$, with $\delta>0$. Then
\[
\abs{f'(r)-g'(r)}\le g'(r+\delta)-g'(r-\delta)+\frac1\delta\sum_{u\in\{r-\delta,r,r+\delta\}}\abs{f(u)-g(u)}.
\]
\end{lemma}

\begin{proof}
By convexity $f'(r)\le\frac{f(r+\delta)-f(r)}\delta$ and $\frac{g(r+\delta)-g(r)}\delta\le g'(r+\delta)$, so $f'(r)-g'(r)\le g'(r+\delta)-g'(r)+\frac1\delta(\abs{f(r+\delta)-g(r+\delta)}+\abs{f(r)-g(r)})$. In the same way $g'(r)-f'(r)\le g'(r)-g'(r-\delta)+\frac1\delta(\abs{f(r)-g(r)}+\abs{f(r-\delta)-g(r-\delta)})$. Since $g'$ is nondecreasing, both are at most the claimed bound.
\end{proof}

\begin{proposition}[Overlap concentration]\label{prop:nishimori:overlap}
For $t\in[0,1]$ and $r>0$ let $V_N(t,r)=\E\gibbs{(R_{1*}-\E\gibbs{R_{1*}}_{t,r})^2}_{t,r}$, and let $M_a=\frac12+(\pi a)^{-1/2}$. For every $N\ge1$, $t\in[0,1]$ and $0<2\delta\le a<b$,
\[
\int_a^bV_N(t,r)\dd r\le\frac{2+4\log\frac ba}N+12M_a\Bigl(\delta+2\sqrt{\frac\delta\pi}\Bigr)+\frac{108(b-a)(\lambda^2+6b)}{N\delta^2}.
\]
For fixed $0<a<b$ the choice $\delta=N^{-2/5}$ makes the right side $O(N^{-1/5})$.
\end{proposition}

\begin{proof}
The plan is to control the thermal and then the disorder fluctuations of the derivative of the free entropy in $r$, and to transfer them to the overlap. Fix $t$, write $\gibbs\cdot$ for $\gibbs\cdot_{t,r}$, $\Phi(r)=\phi_N(t,r)$ and $\hat\Phi(r)=\frac1N\log\cZ_{t,r}$.

\emph{Step 1: thermal fluctuations.} As in the proof of \cref{lem:nishimori:derivatives}, $\hat\Phi'=\gibbs L$ and $\hat\Phi''=N\gibbs{(L-\gibbs L)^2}-\gibbs{\ip zx}/(4Nr^{3/2})$, with $L=R_{1*}+\frac{\ip zx}{2N\sqrt r}-\frac12$, $\E\gibbs L=\Phi'=\frac12\E\gibbs{R_{1*}}\in[0,\frac12]$ and $0\le\E\gibbs{\ip zx}\le\sqrt rN$. Taking expectations, $\E\gibbs{(L-\gibbs L)^2}\le\frac{\Phi''(r)}N+\frac1{4Nr}$, and integrating over $[a,b]$, $\int_a^b\E\gibbs{(L-\gibbs L)^2}\dd r\le\frac1{2N}+\frac{\log(b/a)}{4N}$.

\emph{Step 2: disorder fluctuations.} Since $\abs{\gibbs{\ip zx}}\le\norm z_1$, the function $\tilde\Phi=\hat\Phi-\frac{\sqrt r}N\norm z_1$ is convex, and so is its mean $\bar\Phi=\Phi-\sqrt{2r/\pi}$. For $r\in[a,b]$, \cref{lem:nishimori:convex} with $f=\tilde\Phi$, $g=\bar\Phi$, $\Delta=\tilde\Phi-\bar\Phi$ and $\Lambda(r)=\bar\Phi'(r+\delta)-\bar\Phi'(r-\delta)\ge0$ gives
\[
\abs{\gibbs L-\E\gibbs L}\le\Lambda(r)+\frac1\delta\sum_{u\in\{r-\delta,r,r+\delta\}}\abs{\Delta(u)}+\frac{\abs{\norm z_1-\E\norm z_1}}{2N\sqrt r}.
\]
Given $x^*$, $\log\cZ_{t,u}$ is a function of the standard Gaussians $\sqrt NW_{ij}$ ($i<j$), $\sqrt{N/2}\,W_{ii}$ and $z_i$, with partial derivatives $\frac{\lambda\sqrt t}{\sqrt N}\gibbs{x_ix_j}$, $\frac{\lambda\sqrt t}{\sqrt{2N}}$ and $\sqrt u\gibbs{x_i}$ and squared gradient at most $\frac{\lambda^2t}N\binom N2+\frac{\lambda^2t}2+uN\le(\frac{\lambda^2}2+u)N$. The Gaussian Poincaré inequality (\cref{thm:toolkit:poincare}) and \cref{lem:nishimori:gauge} give $\Var\hat\Phi(u)\le(\frac{\lambda^2}2+u)/N$, and with $\Var\norm z_1\le N$ and $u\le\frac32b$, $\E\Delta(u)^2\le2\Var\hat\Phi(u)+\frac{2u}{N^2}\Var\norm z_1\le(\lambda^2+6b)/N$. Next, $0\le\Lambda\le M_a$, because $\Phi'\in[0,\frac12]$ and $r-\delta\ge\frac a2$; and $\int_a^b\Lambda\dd r\le\delta+2\sqrt{\delta/\pi}$, because the integral equals $[\bar\Phi(b+\delta)-\bar\Phi(b-\delta)]-[\bar\Phi(a+\delta)-\bar\Phi(a-\delta)]$, $\Phi$ is nondecreasing and $\frac12$-Lipschitz, and $\sqrt{a+\delta}-\sqrt{a-\delta}\le\sqrt{2\delta}$. Squaring with $(\alpha_1+\alpha_2+\alpha_3)^2\le3(\alpha_1^2+\alpha_2^2+\alpha_3^2)$, which also gives $(\sum_u\abs{\Delta(u)})^2\le3\sum_u\Delta(u)^2$, using $\Lambda^2\le M_a\Lambda$ and $\Var\norm z_1\le N$, and integrating,
\[
\int_a^b\E(\gibbs L-\E\gibbs L)^2\dd r\le3M_a\Bigl(\delta+2\sqrt{\frac\delta\pi}\Bigr)+\frac{27(b-a)(\lambda^2+6b)}{N\delta^2}+\frac{3\log(b/a)}{4N}.
\]

\emph{Step 3: from $L$ to the overlap.} Write $Q=R_{1*}$, $\bar Q=\E\gibbs Q$ and $\bar L=\E\gibbs L=\frac12\bar Q$. Integration by parts in $z$ gives $\E\gibbs{Q\ip zx}=\sqrt rN\,\E(\gibbs Q-\gibbs{R_{1*}R_{12}})$, so $\E\gibbs{LQ}=\E\gibbs{Q^2}-\frac12\E\gibbs{R_{1*}R_{12}}$, and the Nishimori identity turns this into $\E\gibbs{R_{12}^2}-\frac12\E\gibbs{R_{12}R_{13}}$. Since $V_N=\E\gibbs{R_{12}^2}-\bar Q^2$,
\[
\E\gibbs{(L-\bar L)(Q-\bar Q)}=\frac12V_N+\frac12\E\bigl(\gibbs{R_{12}^2}-\gibbs{R_{12}R_{13}}\bigr)\ge\frac12V_N,
\]
because $N^2(\gibbs{R_{12}^2}-\gibbs{R_{12}R_{13}})=\norm{\mathbf C}_F^2+m^\top\mathbf Cm\ge0$, with $m=\gibbs x$ and the posterior covariance $\mathbf C=\gibbs{xx^\top}-mm^\top$. The Cauchy--Schwarz inequality then gives $V_N\le4\E\gibbs{(L-\bar L)^2}=4\E\gibbs{(L-\gibbs L)^2}+4\E(\gibbs L-\bar L)^2$, and Steps 1 and 2 give the claim. With $\delta=N^{-2/5}$, $\delta\le\sqrt\delta=N^{-1/5}$ and $N\delta^2=N^{1/5}$.
\end{proof}

By the Nishimori identity $R_{12}$ has the law of $R_{1*}$, so $V_N$ is also the variance of the overlap of two replicas: with a side channel, $P(q)$ concentrates on one value for most $r$, at every $\lambda$. This is the rigorous form of replica symmetry.

\subsection{The matching bound}\label{sec:nishimori:matching}

To make Guerra's remainder a variance around the mean, the trial overlap follows the model along the path: the \term{adaptive interpolation}\index{interpolation!adaptive} of \textcite{BarbierMacris2019}, in the continuous form of \textcite{BarbierEtAl2019}.

\begin{proof}[Proof of part 2 of \cref{thm:nishimori:z2}]
The plan is to choose the trial overlap equal to the mean overlap of the interpolating model, so that the remainder in \cref{prop:nishimori:interpolation} is $V_N$, and to average over the side channel so that \cref{prop:nishimori:overlap} applies. Let $\lambda>0$, $N\ge11$, $s=N^{-1/10}$ and $\delta=N^{-2/5}$, so that $2\delta\le s\le1$, and let $G_N(t,r)=\E\gibbs{R_{1*}}_{t,r}\in[0,1]$.

\emph{The adaptive path.} $G_N$ is continuous and nondecreasing in $r$ (\cref{lem:nishimori:derivatives}), and Lipschitz in $r\ge\frac s2$ uniformly in $t$, because $\partial_rG_N=2\partial_r^2\phi_N$ and, by Step 1 of the proof of \cref{prop:nishimori:overlap}, $\abs{\partial_r^2\phi_N}\le N\E\gibbs{L^2}+\frac1{4r}$ with $\abs L\le\frac32+\norm z_1/(2N\sqrt r)$. Extend $G_N(t,\cdot)$ to $r<\frac s2$ by its value at $\frac s2$. For $\epsilon\in[s,2s]$ the Cauchy--Lipschitz theorem gives a unique continuously differentiable solution $r_\epsilon$ of $r'(t)=-\lambda^2G_N(t,r(t))$ on $[0,1]$ with $r(1)=\epsilon$; it depends continuously on $(t,\epsilon)$, and $\epsilon\le r_\epsilon\le\epsilon+\lambda^2$. It expands distances: for $\epsilon_1<\epsilon_2$, $\Delta=r_{\epsilon_2}-r_{\epsilon_1}$ satisfies
\[
\Delta(t)=\epsilon_2-\epsilon_1+\lambda^2\int_t^1\bigl[G_N(u,r_{\epsilon_2}(u))-G_N(u,r_{\epsilon_1}(u))\bigr]\dd u,
\]
so if $t_0$ is the infimum of the $t$ with $\Delta>0$ on $[t,1]$, then $\Delta\ge\epsilon_2-\epsilon_1$ on $(t_0,1]$, hence at $t_0$ by continuity, and $t_0=0$, since otherwise $\Delta>0$ slightly to the left of $t_0$.

\emph{The sum rule.} With $q(t)=G_N(t,r_\epsilon(t))$, $r_\epsilon(t)=\epsilon+\lambda^2\int_t^1q$, and the remainder in \cref{prop:nishimori:interpolation} is $\frac{\lambda^2}4\int_0^1V_N(t,r_\epsilon(t))\dd t$. Three inequalities simplify the rest: $\phi_N\le\phi_N(1,\epsilon)$ because $\partial_r\phi_N\ge0$; $\int_0^1q^2\ge\bar q^2$ for $\bar q=\int_0^1q\in[0,1]$; and $\psi(\epsilon+\lambda^2\bar q)\le\psi(\lambda^2\bar q)+\frac\epsilon2$ because $0\le\psi'\le\frac12$. Since $\frac\epsilon2\le s$,
\[
\phi_N\le\phi_{\mathrm{RS}}(\lambda,\bar q)+s+\frac{\lambda^2}4\int_0^1V_N(t,r_\epsilon(t))\dd t.
\]

\emph{Averaging over the side channel.} For fixed $t$ the map $\rho(\epsilon)=r_\epsilon(t)$ is continuous and expanding, and $V_N(t,\cdot)\ge0$ is continuous, so comparing Riemann sums, $\sum_kV_N(t,\rho(\epsilon_k))(\epsilon_{k+1}-\epsilon_k)\le\sum_kV_N(t,\rho(\epsilon_k))(\rho(\epsilon_{k+1})-\rho(\epsilon_k))$, gives $\int_s^{2s}V_N(t,r_\epsilon(t))\dd\epsilon\le\int_s^{2s+\lambda^2}V_N(t,r)\dd r$. Averaging the last bound over $\epsilon\in[s,2s]$ and applying \cref{prop:nishimori:overlap} with $a=s$ and $b=2s+\lambda^2$,
\begin{align*}
\phi_N-\max_q\phi_{\mathrm{RS}}(\lambda,q)\le s+\frac{\lambda^2}{4s}\Bigl[&\frac{2+4\log(2+\lambda^2/s)}N+12M_s\Bigl(\delta+2\sqrt{\frac\delta\pi}\Bigr)\\
&+\frac{108(s+\lambda^2)(7\lambda^2+12s)}{N\delta^2}\Bigr].
\end{align*}
With $s=N^{-1/10}$ and $\delta=N^{-2/5}$ the three terms in the bracket contribute $O(N^{-9/10}\log N)$, $O(s^{-3/2}\sqrt\delta)=O(N^{-1/20})$ and $O(s^{-1}N^{-1/5})=O(N^{-1/10})$, with constants depending only on $\lambda$. With part 1 this proves part 2.
\end{proof}

The rate $N^{-1/20}$ is an artifact of the argument; \cref{fig:nishimori:measured}(b) suggests a gap of order $1/N$.

\subsection{Uniqueness and the minimum mean-square error}\label{sec:nishimori:mmse}

For signs the maximizer is unique for every $\lambda$, because the binary mmse is a mixture of decaying exponentials.\index{complete monotonicity}

\begin{lemma}[The binary mmse is completely monotone]\label{lem:nishimori:cm}
For every $\gamma\ge0$, with the probability measure $\mu(\dd\omega)=\frac12\operatorname{sech}(\frac{\pi\omega}2)\dd\omega$,
\[
1-G(\gamma)=e^{-\gamma/2}\,\E\operatorname{sech}(\sqrt\gamma Z)=\int_\R e^{-(1+\omega^2)\gamma/2}\,\mu(\dd\omega).
\]
Consequently $G$ is strictly concave on $[0,\infty)$, with $G(0)=0$ and $G'(0)=1$.
\end{lemma}

\begin{proof}
The plan is to remove the drift by a change of measure and to expand $\operatorname{sech}$ in Fourier modes. The density of $\gamma+\sqrt\gamma Z$ at $h$ is $e^{h-\gamma/2}$ times that of $\sqrt\gamma Z$, so $\E f(\gamma+\sqrt\gamma Z)=e^{-\gamma/2}\E[e^{\sqrt\gamma Z}f(\sqrt\gamma Z)]$; with $f=1-\tanh$ and $(1-\tanh u)e^u=\operatorname{sech}u$ this is the first identity. The function $\operatorname{sech}$ is its own Fourier transform up to scaling, $\int_\R\operatorname{sech}(u)e^{-i\omega u}\dd u=\pi\operatorname{sech}(\pi\omega/2)$ (a residue computation), so $\operatorname{sech}u=\int\cos(\omega u)\mu(\dd\omega)$, and Fubini's theorem with $\E\cos(\omega\sqrt\gamma Z)=e^{-\omega^2\gamma/2}$ gives the second. Differentiating under the integral, which the exponential decay of $\mu$ allows, $G''(\gamma)=-\frac14\int(1+\omega^2)^2e^{-(1+\omega^2)\gamma/2}\mu(\dd\omega)<0$, and $G'(0)=\frac12\int(1+\omega^2)\mu(\dd\omega)=1$, because $\int\omega^2\mu(\dd\omega)=-\operatorname{sech}''(0)=1$.
\end{proof}

\begin{proof}[Proof of parts 3 and 4 of \cref{thm:nishimori:z2}]
The plan is to read the maximizer off the concavity of $G$, and to obtain the error by differentiating the limit of the free entropy in $\theta=\lambda^2$.

\emph{Part 3.} Since $\psi'=\frac12G$, $\partial_q\phi_{\mathrm{RS}}=\frac{\lambda^2}2h(q)$ with $h(q)=G(\lambda^2q)-q$, which is strictly concave, with $h(0)=0$, $h'(0)=\lambda^2-1$ and $h(1)<0$. If $\lambda\le1$, then $h<0$ on $(0,1]$ and $q_*=0$. If $\lambda>1$, then $h>0$ near $0$, and a strictly concave function vanishing at $0$ has at most one other zero; so $h$ has exactly one zero $q_\circ\in(0,1)$, is positive before it and negative after it, and $q_*=q_\circ$. The zeros of $h$ in $(0,1]$ are the solutions of \eqref{eq:nishimori:fixedpoint}.

\emph{Part 4.} At $\lambda=0$ the posterior is uniform and $\mathrm{MMSE}_N=1-\frac1N$. For $\theta=\lambda^2>0$, \cref{lem:nishimori:derivatives} at $r=0$ gives $\frac{\dd\phi_N}{\dd\theta}=\frac14\E\gibbs{R_{1*}^2}=\frac14(1-\mathrm{MMSE}_N)$, nondecreasing in $\theta$, so $\phi_N$ is convex in $\theta$. Its limit $\Phi(\theta)=\max_{q\in[0,1]}\tilde\phi(\theta,q)$, with $\tilde\phi(\theta,q)=\psi(\theta q)-\frac{\theta q^2}4$, has $\Phi'(\theta)=\frac14q_*^2$: the difference $\Phi(\theta')-\Phi(\theta)$ lies between $\tilde\phi(\theta',q)-\tilde\phi(\theta,q)$ at $q=q_*(\theta)$ and at $q=q_*(\theta')$, both are $(\theta'-\theta)\,\partial_\theta\tilde\phi$ at intermediate points, $q_*(\theta')\to q_*(\theta)$ by uniqueness and compactness, and $\partial_\theta\tilde\phi=q\psi'(\theta q)-\frac{q^2}4=\frac{q^2}4$ at $q=q_*$ by the fixed-point equation. Finally, if convex $f_N$ converge pointwise to $f$, then $\frac{f(\theta)-f(\theta-u)}u\le\liminf f_N'(\theta)\le\limsup f_N'(\theta)\le\frac{f(\theta+u)-f(\theta)}u$ for small $u>0$, and both bounds tend to $f'(\theta)$. Hence $\E\gibbs{R_{1*}^2}\to q_*^2$.
\end{proof}

\subsection{What the general proof adds}

Signs gave us a constant norm $\norm x^2=N$, bounded spins and the gauge symmetry. For priors with bounded support the sum rule holds verbatim (\cref{ex:nishimori:general}), and \textcite{BarbierMacris2019} control the variance and the self-overlaps with a Gaussian side channel, interpolating in $K$ small steps. \Textcite{LelargeMiolane2019} instead reveal each $x^*_i$ with a small probability and run the Aizenman--Sims--Starr cavity computation \parencite{AizenmanSimsStarr2003}, comparing sizes $N$ and $N+1$, for priors with finite support, which approximate every prior with a finite second moment; and they remove the perturbation: for bounded priors, at points of differentiability, $\E\gibbs{(R_{1*}^2-q_*^2)^2}\to0$ (their Theorem~2).

\section{Measuring it}\label{sec:nishimori:measuring}

The script \texttt{04-nishimori-exact.py} enumerates the posterior at $N=6$, $8$, $10$, $12$: for ten values of $\lambda$ it draws $1200$ instances and records the free entropy, the matrix error and both sides of $\E\gibbs{R_{1*}^2}=\E\gibbs{R_{12}^2}$. Bayes' rule writes $\log\cZ_N=B-\log(2^N\Prob(x^*\mid A))$ with $B=H_N(x^*)-\frac{\lambda^2N}4$ of known mean and variance $\frac{\lambda^2N}2$, and $B$ serves as a control variate that shrinks the intervals by a factor of up to $19$. The script \texttt{04-nishimori-amp.py} runs Bayes-optimal message passing (\cref{ch:amp}) from a spectral start on ten instances at $N=2000$ per value of $\lambda$, and records the matrix error of its estimate after $40$ iterations and that of principal components, $(1-\lambda^{-2})_+Nv_1v_1^\top$; message passing reaches the Bayes error for $\lambda>1$ \parencite{DeshpandeAbbeMontanari2016,MontanariVenkataramanan2021}.

\begin{figure}[t]
\begin{fullwidth}
\centering
\begin{tikzpicture}
\begin{groupplot}[group style={group size=2 by 1,horizontal sep=17mm},rplot,width=80mm,height=52mm]
\nextgroupplot[xlabel={$\lambda$},ylabel={matrix mean-square error},xmin=0,xmax=3.08,ymin=-0.03,ymax=1.06,
  title={(a) the Bayes error},legend pos=south west,legend style={font=\scriptsize}]
\addplot[physpred] table[x=lam,y=mmse] {figures/data/04-nishimori-curves.dat};
\addplot[proofpred] table[x=lam,y=spec] {figures/data/04-nishimori-curves.dat};
\addplot[measured,ebars] table[x=lam,y=amp,y error plus=ampci,y error minus=ampci] {figures/data/04-nishimori-amp.dat};
\addplot[measured3,ebars] table[x=lam,y=spec,y error plus=specci,y error minus=specci] {figures/data/04-nishimori-amp.dat};
\addplot[measured2,ebars] table[x=lam,y=mmse12,y error plus=mmse12ci,y error minus=mmse12ci] {figures/data/04-nishimori-exact.dat};
\legend{prediction $1-q_*^2$ (proved),PCA limit (proved),AMP ($N=2000$),PCA ($N=2000$),exact posterior ($N=12$)}
\nextgroupplot[xlabel={$\lambda$},ylabel={$N\bigl(\phi_N-\max_q\phi_{\mathrm{RS}}(\lambda,q)\bigr)$},xmin=0,xmax=3.08,ymin=-0.05,ymax=1.25,
  title={(b) the gap of the free entropy},legend pos=south east,legend style={font=\scriptsize}]
\addplot[proofpred] table[x=lam,y=bound] {figures/data/04-nishimori-below.dat};
\addplot[proofpred,forget plot] coordinates {(0,0) (3.08,0)};
\addplot[physpred2] table[x=lam,y=heur] {figures/data/04-nishimori-below.dat};
\addplot[measured3,ebars] table[x=lam,y=gap8,y error plus=gap8ci,y error minus=gap8ci] {figures/data/04-nishimori-exact.dat};
\addplot[measured2,ebars] table[x=lam,y=gap10,y error plus=gap10ci,y error minus=gap10ci] {figures/data/04-nishimori-exact.dat};
\addplot[measured,ebars] table[x=lam,y=gap12,y error plus=gap12ci,y error minus=gap12ci] {figures/data/04-nishimori-exact.dat};
\legend{proved bounds,lognormal heuristic,$N=8$,$N=10$,$N=12$}
\end{groupplot}
\end{tikzpicture}
\end{fullwidth}
\caption{(a) Matrix mean-square error. Carnelian: $1-q_*^2$ (\cref{thm:nishimori:z2}); lapis dashed: the limit $1-(1-\lambda^{-2})_+^2$ of principal components (\cref{thm:spiked:bbp}), an upper bound on the Bayes error. Dots: message passing at $N=2000$, ten instances per point (seeds $4000+100k+r$); squares: principal components on the same instances; circles: exact Bayes error at $N=12$, $1200$ instances per point (seeds $1000N+k$). (b) $N$ times the gap of the free entropy, by exact enumeration ($1200$ instances per point, control variate). Lapis dashed: the proved bounds $0$ and, for $\lambda<1$, $-\frac12\log(1-\lambda^2)$ (\cref{prop:nishimori:below}); carnelian dotted: the lognormal heuristic $-\frac14\log(1-\lambda^2)$. Student-$t$ $95\%$ intervals. Scripts \texttt{04-nishimori-exact.py}, \texttt{-amp.py}, \texttt{-rs.py}.}\label{fig:nishimori:measured}
\end{figure}

In panel (a), message passing follows the prediction: ten of its eleven intervals contain $1-q_*^2$, and the eleventh, at $\lambda=2$, misses it by $0.001$ ($0.172\pm0.011$ against $0.160$). Above the threshold the Bayes error is far below that of principal components, $0.160$ against $0.4375$ at $\lambda=2$, and all eleven spectral intervals contain their proved limit; near the threshold the two curves touch (\cref{ex:nishimori:threshold}). The exact error at $N=12$ differs from its limit by up to $0.23$, at the threshold ($0.772$ against $1$); below the threshold it is at most $1-\frac1N=0.917$, because the posterior mean knows the diagonal of $x^*x^{*\top}$, a twelfth of the matrix, and the measured $0.85$ to $0.91$ approach this as $\lambda$ decreases.

In panel (b), all forty values of $N(\phi_N-\max_q\phi_{\mathrm{RS}})$ are positive, as part 1 of \cref{thm:nishimori:z2} requires, and of order one: $0.65$ to $1.08$ for $\lambda\ge1.25$. Below the threshold they lie under the bound of \cref{prop:nishimori:below} and near half of it ($0.011$--$0.020$, $0.053$--$0.074$ and $0.170$--$0.203$ at $\lambda=0.25$, $0.5$, $0.75$, against $0.016$, $0.072$, $0.207$), which is the mean under the spiked law of a Gaussian log-likelihood ratio with second moment $(1-\lambda^2)^{-1/2}$; \textcite{ElAlaouiKrzakalaJordan2020} prove the Gaussian limit below the threshold, and the convergence of the moments is the heuristic step. The Nishimori identity holds within the noise: the forty differences are below $0.009$ in absolute value except one, $-0.016\pm0.009$ at $N=12$, $\lambda=1.75$.\sidenote{The per-instance differences are skewed: at that point their median is $+0.04$ and their $0.1\%$ quantile between $-0.7$ and $-0.6$, so a few instances on which the posterior is confident and wrong carry the mean, and Student intervals undercover. Two fresh sets of $3000$ instances give $+0.0006\pm0.0048$ and $-0.0001\pm0.0047$.} The scripts run in about $40$, $70$ and $10$ seconds, and no measurement contradicts a theorem of the chapter.

\section{The status board}\label{sec:nishimori:status}

As of 2026 everything in the prediction is a theorem, for signs and for every prior with a finite second moment. What remains physics is the principle that the Nishimori line never breaks replica symmetry, and the hardness of the hard phase; off the line the replica-symmetric answer is wrong at low temperature.

\begin{statusboard}
\statusrow{Nishimori identity, at every $N$}{\proved}{Bayes' rule (\cref{lem:nishimori:nishimori}); gauge form: \textcite{Nishimori1981}}
\statusrow{$P(q)=P(m)$ on the Nishimori line}{\proved}{the identity with $f=\ind\{R_{12}\in B\}$; $\pm J$ model: \textcite{NishimoriSherrington2001}}
\statusrow{No replica-symmetry breaking on the Nishimori line, in every model}{\conjectured}{\cref{sec:nishimori:physics}; open in general \parencite[\S II.G]{ZdeborovaKrzakala2016}}
\statusrow{Mutual information and error of $\Z_2$ synchronization}{\proved}{\textcite{DeshpandeAbbeMontanari2016}; \textcite{LelargeMiolane2019}; \cref{thm:nishimori:z2}}
\statusrow{The same for every prior with finite second moment}{\proved}{\textcite[Thm.~1]{LelargeMiolane2019}; bounded support: \textcite{BarbierMacris2019}}
\statusrow{Continuous transition at $\lambda=1$, with a unique maximizer}{\proved}{from the formula; uniqueness: \cref{lem:nishimori:cm} (this book)}
\statusrow{Overlap concentration (rigorous replica symmetry)}{\proved}{side channel: \cref{prop:nishimori:overlap}, \textcite{Barbier2021}; without: \textcite[Thm.~2]{LelargeMiolane2019}}
\statusrow{Replica-symmetric formula for a posterior at the wrong temperature}{\knownfalse}{$\lambda=0$, $\beta>1$ \parencite{Toninelli2002}; correct formula: \textcite{Chen2014}}
\statusrow{The continuation $n\to0$ and the exchange of limits}{\heuristicstatus}{no counterpart: the proof never continues in $n$}
\statusrow{Sparse priors: the hard phase is hard for efficient algorithms}{\conjectured}{\textcite[\S IV C]{LesieurKrzakalaZdeborova2017}; evidence in \cref{ch:spiked}}
\end{statusboard}

\section{Where the translation fails}\label{sec:nishimori:fails}

\paragraph{What was assumed.} Bayes optimality does the work: the statistician knows the prior, the signal-to-noise ratio and the noise, and computes the exact posterior. The prediction assumes it too, and every theorem here uses it through the Nishimori identity. The model is idealized as usual: a product prior, every pair observed, Gaussian noise, rank one, $N\to\infty$ at fixed $\lambda$. Finite rank is treated by \textcite[\S2.8]{LelargeMiolane2019}, and other noise channels reduce to the Gaussian one by channel universality, proved by \textcite{KrzakalaXuZdeborova2016} for priors with finite support and channels with bounded second and third log-derivatives. Our proof also assumed signs.

\paragraph{Where the physics is wrong, or not yet right.} Off the Nishimori line the argument collapses: a posterior at the wrong temperature is a planted Sherrington--Kirkpatrick model, the Nishimori identity fails, the remainder loses its sign (\cref{ex:nishimori:sk}), and at $\lambda=0$, $\beta>1$ the replica-symmetric formula is false \parencite{Toninelli2002}; the free entropy is then the formula of \textcite{Chen2014}. On the line, replica symmetry is a theorem only for dense models with independent priors \parencite[\S II.G]{ZdeborovaKrzakala2016}. ``Self-averaging'' is literally false without a perturbation: for signs $\gibbs{R_{1*}}=0$ in every instance, while $\abs{R_{1*}}$ concentrates; physicists break the symmetry by hand, the proof with the side channel. Replica-symmetric statics says nothing about how the posterior is reached: for sparse priors a second maximum can trap an efficient algorithm. And at accessible sizes the corrections are large (up to $0.23$ in the error at $N=12$), and the replica computation is silent about them.

\paragraph{Open problems.} Three are posed in the literature. Whether replica symmetry holds on the Nishimori line of finite-dimensional models \parencite[\S II.G]{ZdeborovaKrzakala2016}. Whether a polynomial-time algorithm can beat message passing in the hard phase of sparse priors, which \textcite[\S IV C]{LesieurKrzakalaZdeborova2017} conjecture to be impossible (\cref{ch:spiked}). And whether, for $\Z_2$ synchronization, the local minimizer of the TAP free energy at the posterior mean, which \textcite{CelentanoFanMei2023} construct for every $\lambda>1$, is global, as they conjecture from numerics (\cref{ch:amp}).

\begin{exercises}
\item \consolidation (The scalar channel.) Let $h=\gamma+\sqrt\gamma Z$, $\gamma>0$. (a) Show that the density $p$ of $h$ satisfies $p(-h)=e^{-2h}p(h)$, and deduce $\E[f(h)\tanh h]=\E[f(h)\tanh^2h]$ for every bounded even $f$; in particular $G(\gamma)=\E\tanh^2h$. (b) Prove $\psi'=\frac12G$ by integration by parts. (c) Check that $\phi_N(0,r)=\psi(r)$.\label{ex:nishimori:scalar}

\item \consolidation (Information and error.) (a) Prove \eqref{eq:nishimori:mi} and the second form of \eqref{eq:nishimori:mmse}. (b) Show that $\frac{\dd}{\dd\lambda^2}\frac1NI(x^*;A)=\frac14\mathrm{MMSE}_N$, and check numerically that $\frac14\int_0^{2.25}(1-q_*(\sqrt\theta)^2)\dd\theta=0.4942$, the limit of $\frac1NI$ at $\lambda=1.5$.\label{ex:nishimori:immse}

\item \consolidation (The Gaussian prior.) For $P_0=\Normal(0,1)$ show that $\psi_{P_0}(\gamma)=\frac12(\gamma-\log(1+\gamma))$, $1-\mathrm{mmse}(\gamma)=\frac\gamma{1+\gamma}$ and $q_*=(1-\lambda^{-2})_+$, so that principal components are Bayes-optimal, and that the right side of \cref{thm:nishimori:rsformula} is $\frac{\lambda^2}4$ for $\lambda\le1$ and $\log\lambda+\frac1{4\lambda^2}$ for $\lambda>1$.\label{ex:nishimori:gaussian}

\item \consolidation (Near the threshold.) From $G(\gamma)=\gamma-\gamma^2+\frac53\gamma^3+O(\gamma^4)$, show that $q_*=2\epsilon-\frac13\epsilon^2+O(\epsilon^3)$ at $\lambda=1+\epsilon$, and that $1-q_*^2=1-4\epsilon^2+\frac43\epsilon^3+O(\epsilon^4)$, while principal components give $1-4\epsilon^2+12\epsilon^3+O(\epsilon^4)$.\label{ex:nishimori:threshold}

\item \byhand (The series of $G$.) (a) Show that the moments $m_{2k}=\int\omega^{2k}\mu(\dd\omega)$ of the measure of \cref{lem:nishimori:cm} are $(-1)^k\operatorname{sech}^{(2k)}(0)$, the Euler numbers $\abs{E_{2k}}=1,1,5,61,1385,\dots$ (b) Deduce that, as $\gamma\to0$, $G$ has the asymptotic expansion $\sum_{k\ge1}(-1)^{k+1}c_k\gamma^k$, where $c_k=\frac1{2^kk!}\sum_{j\le k}\binom kj\abs{E_{2j}}$, and check the coefficients $1$, $-1$, $\frac53$, $-\frac{13}3$. (c) Show that the series diverges for every $\gamma>0$. \hint{$\abs{E_{2k}}\ge(2k)!\,(2/\pi)^{2k+1}$.}\label{ex:nishimori:series}

\item \byhand (Replica symmetry is stable on the Nishimori line.) (a) Using \cref{ex:nishimori:scalar}, show that $G'(\gamma)=\E\operatorname{sech}^4(\gamma+\sqrt\gamma Z)$. (b) The de Almeida--Thouless condition of \cref{ch:sk}, with the ferromagnetic field $\beta J_0m$ in place of the external field, is $\beta^2\E\operatorname{sech}^4(\beta\sqrt qZ+\beta J_0m)\le1$. Show that on the Nishimori line, with $m=q=q_*$, it reads $\lambda^2G'(\lambda^2q_*)\le1$, and deduce from \cref{lem:nishimori:cm} that it holds for every $\lambda$, strictly for $\lambda\ne1$ ($0.838$ at $\lambda=1.1$, $0.457$ at $\lambda=1.5$).\label{ex:nishimori:at}

\item \byhand (Every prior.) Let $P_0$ have bounded support and $\E X^2=1$. (a) Show that \eqref{eq:nishimori:replicated} holds verbatim with $x^0,\dots,x^n$ independent with law $P_0^{\otimes N}$: the terms in $\norm{x^a}^4$ cancel against the diagonal of the GOE. (b) Redo the replica-symmetric computation and obtain $\psi_{P_0}$. (c) With the evidence of the interpolating model as in \cref{sec:nishimori:interp} but averaged over $P_0^{\otimes N}$, check \cref{lem:nishimori:derivatives} and the sum rule with $\psi_{P_0}$ for $\psi$, and deduce Guerra's bound at every $N$.\label{ex:nishimori:general}

\item \translate (Off the Nishimori line.) Let $\lambda=0$, so that $A=W$, and take the Gibbs measure proportional to $e^{\frac\beta2x^\top Wx}$, the posterior computed from pure noise by a statistician who believes the ratio is $\beta$. With $H_t(x)=\frac{\beta\sqrt t}2x^\top Wx+\beta\sqrt{q(1-t)}\ip zx$, show that $\frac{\dd}{\dd t}\frac1N\E\log\sum_xe^{H_t(x)}=\frac{\beta^2}4((1-q)^2-\E\gibbs{(R_{12}-q)^2}_t)$, deduce Guerra's upper bound (\cref{thm:sk:rsbound}) at zero field, and find the step of the proof of \cref{lem:nishimori:derivatives} that fails.\label{ex:nishimori:sk}

\item \translate (Self-averaging of a derivative.) Let $f_N$ be random convex functions on an open interval $I$ with $\sup_I\Var f_N\le c/N$, and $g_N=\E f_N$ with $\abs{g_N'}\le1$. (a) Show with \cref{lem:nishimori:convex} that $\int_J\E\abs{f_N'-g_N'}=O(N^{-1/4})$ on compact intervals $J\subset I$. (b) Show that the integral cannot be dropped, with $f_N(u)=\abs{u-U_N}$ and $U_N$ uniform on $[u_0-\frac1N,u_0+\frac1N]$.\label{ex:nishimori:selfavg}

\item \measurement (Overlap distributions.) Modify the enumeration script of \cref{sec:nishimori:measuring} to record at $N=12$, for $\lambda=0.5$ and $2$, the laws of $R_{12}$ and $R_{1*}$ under $\E\gibbs\cdot$. Check that they agree ($P(q)=P(m)$), that at $\lambda=0.5$ both are one bump around $0$ of width about $N^{-1/2}$, and that at $\lambda=2$ most of their mass sits at $\pm1$ and $\pm\frac56$. Add a side channel with $r=0.5$: what happens to the symmetry between $R$ and $-R$, and why?\label{ex:nishimori:histograms}

\item \research (The mismatched posterior.) A statistician who believes the ratio is $\beta$ when it is $\lambda$ computes the Gibbs measure proportional to $e^{\frac\beta2x^\top Ax}$. (a) Show that after the gauge transformation it is the Sherrington--Kirkpatrick model at inverse temperature $\beta$ plus $\frac{\beta\lambda}N\sum_{i<j}\sigma_i\sigma_j$. (b) Show that the replica-symmetric equations $m=\E\tanh(\beta\sqrt qZ+\beta\lambda m)$ and $q=\E\tanh^2(\beta\sqrt qZ+\beta\lambda m)$ have the solution $m=q=q_*$ when $\beta=\lambda$. (c) Map where this solution violates the de Almeida--Thouless condition, compare with the formula of \textcite{Chen2014}, and find where the mismatched posterior mean does worse than principal components.\label{ex:nishimori:mismatch}
\end{exercises}

\notesandrefs
\textcite{Nishimori1981} introduced the gauge argument for the $\pm J$ spin glass and computed the internal energy exactly on the line that now carries his name; his book \parencite{Nishimori2001} develops the subject. \textcite{Iba1999} identified the line with Bayes-optimal inference, and the review of \textcite{ZdeborovaKrzakala2016} builds the statistical physics of inference on it.

For signs, the gauge-transformed model is the Sherrington--Kirkpatrick model with a ferromagnetic mean of the couplings, which \textcite{SherringtonKirkpatrick1975} already allowed; on the Nishimori line its replica-symmetric equations are solved by $m=q=q_*$ (\cref{ex:nishimori:mismatch}). For general priors the formula was derived by \textcite{LesieurKrzakalaZdeborova2017} and in their earlier work. \Textcite{BarbierMacris2019} credit \textcite{KoradaMacris2009}, who solved the gauge-symmetric $p$-spin model on the complete graph, with a full proof of the binary case. \Textcite{DeshpandeAbbeMontanari2016} proved the formula for signs by message passing and interpolation; \textcite{BarbierEtAl2016}, for discrete priors with at most three stationary points of the potential, by spatial coupling; \textcite{KrzakalaXuZdeborova2016} proved Guerra's bound at every $N$; \textcite{LelargeMiolane2019} proved the formula for every prior with a finite second moment; \textcite{BarbierMacris2019} introduced the adaptive interpolation, which \textcite{BarbierEtAl2019} carried, in continuous form, to generalized linear models (\cref{ch:glm}); and \textcite{Barbier2021} proved the concentration of the overlap matrix in optimal Bayesian inference. The interpolation goes back to \textcite{GuerraToninelli2002} and \textcite{Guerra2003}; below the threshold even detection is impossible \parencite{ElAlaouiKrzakalaJordan2020}.

The proof of \cref{thm:nishimori:z2} assembles standard pieces, following \textcite{BarbierMacris2019} and \textcite{BarbierEtAl2019}; the gauge argument for the variance, the constants of \cref{prop:nishimori:overlap}, \cref{prop:nishimori:below} and the uniqueness argument through \cref{lem:nishimori:cm} are the book's own. The representation of the binary mmse as a mixture of exponentials is elementary, but we did not find it in the sources we read.

\paragraph{Translating the literature.} \Textcite{LelargeMiolane2019} observe $Y_{ij}=\sqrt{\lambda/n}\,X_iX_j+Z_{ij}$, $i<j$: their $\lambda$ is our $\lambda^2$, their $n$, $X$ and $x\cdot X$ are our $N$, $x^*$ and $R_{1*}$, their $F(\lambda,q)$ is our $\phi_{\mathrm{RS}}(\sqrt\lambda,q)$, and their free energy $F_n=\frac1n\E\log Z_n$ is our free entropy up to $O(1/N)$ from the diagonal. \Textcite{BarbierMacris2019} write $\Delta=1/\lambda^2$ for the noise variance, $f_n=-\frac1n\E\ln\cZ$, and $f_{\mathrm{RS}}(m;\Delta)=-\phi_{\mathrm{RS}}(\lambda,m)$, so their minimum over $m$ is our maximum over $q$. \Textcite{DeshpandeAbbeMontanari2016} write $\lambda$ for our $\lambda^2$, keep $\gamma=\lambda^2q$, and write the limit of $\frac1NI$ as $\Psi(\gamma_*,\lambda)=\frac\lambda4+\frac{\gamma_*^2}{4\lambda}-\frac{\gamma_*}2+I(\gamma_*)$ with $I(\gamma)=\frac\gamma2-\psi(\gamma)$. \Textcite{LesieurKrzakalaZdeborova2017} have $\Delta/\rho^2=1/\lambda^2$, with $\rho=\E X^2$; \textcite{MontanariVenkataramanan2021} use our $\lambda$. Physicists' $P(q)$ and $P(m)$ are the laws of $R_{12}$ and $R_{1*}$ under $\E\gibbs\cdot$.

\printchapterbib
\end{refsection}
